% ============================================================
%  SIAM Journal Article Skeleton
%  Document class: siamart220329  (v1.4.4, 2022/03/29)
%  Usage: pdflatex main && bibtex main && pdflatex main && pdflatex main
% ============================================================
\documentclass[review,hidelinks,onefignum,onetabnum]{siamart220329}

% ----------------------------------------------------------
% Standard packages (subset – add more as needed)
% ----------------------------------------------------------
\usepackage{amsfonts}          % \mathbb, \mathcal
\usepackage{amsopn}            % \DeclareMathOperator
\usepackage{graphicx}          % \includegraphics
\usepackage{algorithmic}       % algorithmic environment

% Serial/Oxford comma for cleveref lists
\newcommand{\creflastconjunction}{, and~}

% Custom theorem/remark environments (optional, add as needed)
\newsiamremark{remark}{Remark}
\newsiamthm{claim}{Claim}

% ----------------------------------------------------------
% Running headers  (short title | short author list)
% ----------------------------------------------------------
\headers{SHORT TITLE HERE}{A. AUTHOR, B. COAUTHOR, AND C. THIRDAUTHOR}

% ----------------------------------------------------------
% Title and authors
% ----------------------------------------------------------
\title{%
  Full Paper Title Goes Here%
  \thanks{%
    Submitted to the editors \today.%
    \funding{%
      This work was supported by [FUNDING AGENCY] under grant [NUMBER].%
    }%
  }%
}

\author{%
  Alice Author%
    \thanks{%
      Department of Mathematics, University of Somewhere, City, State~00000
      (\email{aauthor@somewhere.edu}).%
    }%
  \and
  Bob Coauthor%
    \thanks{%
      Institute for Applied Science, Another University, City, State~00000
      (\email{bcoauthor@another.edu}).%
    }%
  \and
  Carol Thirdauthor%
    \footnotemark[3]%   % shares address with Bob (same \thanks)
}

% Optional: set PDF metadata
\ifpdf
\hypersetup{
  pdftitle={Full Paper Title Goes Here},
  pdfauthor={A. Author, B. Coauthor, and C. Thirdauthor}
}
\fi

% ----------------------------------------------------------
% Cross-reference to supplement (comment out if no supplement)
% ----------------------------------------------------------
% \externaldocument[][nocite]{supplement}

% ===========================================================
\begin{document}
\maketitle

% ----------------------------------------------------------
% Front matter (all three are REQUIRED by SIAM)
% ----------------------------------------------------------
\begin{abstract}
  [PLACEHOLDER: Replace with a self-contained abstract of 150--250 words.
  Abstracts may not cite references, contain displayed equations, or refer
  to section numbers.  Inline math is permitted.]
\end{abstract}

\begin{keywords}
  keyword one, keyword two, keyword three
\end{keywords}

\begin{MSCcodes}
  65F10, 65N22, 90C25
\end{MSCcodes}

% ===========================================================
\section{Introduction}
\label{sec:intro}

[PLACEHOLDER: Motivate the problem, survey related work with citations
\cite{placeholder1}, and state the main contributions clearly.
The paper is organized as follows: \cref{sec:methods} describes \ldots]

% ===========================================================
\section{Methods}
\label{sec:methods}

[PLACEHOLDER: Describe the mathematical framework and algorithms.]

\begin{theorem}[Main result]\label{thm:main}
  Let $A \in \mathbb{R}^{n \times n}$ be symmetric positive definite.
  Then \ldots
\end{theorem}

\begin{proof}
  [PLACEHOLDER: Proof of \cref{thm:main}.]
\end{proof}

% ===========================================================
\section{Numerical experiments}
\label{sec:experiments}

[PLACEHOLDER: Describe experimental setup, datasets, and metrics.
\Cref{tab:results} summarizes the results.]

\begin{table}[htbp]
  \footnotesize
  \caption{Summary of numerical results.}
  \label{tab:results}
  \begin{center}
    \begin{tabular}{|l|c|c|c|} \hline
      Method & Error & Iterations & Time (s) \\ \hline
      Method A & 1.2e-6 & 42 & 0.31 \\
      Method B & 3.5e-7 & 19 & 0.14 \\ \hline
    \end{tabular}
  \end{center}
\end{table}

% ===========================================================
\section{Conclusions}
\label{sec:conclusions}

[PLACEHOLDER: Summarize findings and outline directions for future work.]

% ----------------------------------------------------------
% Appendix (optional -- comment out if unused)
% ----------------------------------------------------------
\appendix
\section{Proof of \cref{thm:main}}
\label{sec:proof}

[PLACEHOLDER: Detailed proof.]

% ----------------------------------------------------------
% Acknowledgments
% ----------------------------------------------------------
\section*{Acknowledgments}

[PLACEHOLDER: Thank colleagues, data providers, and referees.]

% ----------------------------------------------------------
% Bibliography
% ----------------------------------------------------------
\bibliographystyle{siamplain}
\bibliography{references}

\end{document}
